% Faithful official CCC2022 article16 Section4.5 and AppendixB fine-grained main results.
\section*{4.5 Fine-grained case}
Recall the quasilinear-time complexity classes $\mathrm{QL}=\mathrm{DTIME}[\widetilde O(n)]$, $\mathrm{NQL}=\mathrm{NTIME}[\widetilde O(n)]$, and the quasilinear-time analog $\mathrm{QLH}=\bigcup_{\ell\geq0}\Sigma_\ell^{\mathrm{QL}}$ of the polynomial-time hierarchy PH; see, e.g., [36] for more details. Define $\mathrm{BPQL}=\mathrm{BPTIME}[\widetilde O(n)]$, the quasilinear-time version of BPP. Also, define $\mathrm{QLSamp}$ to be the class of distribution families that are quasilinear-time samplable. We show better probabilistic worst-case to average-case reductions under fine-grained average-case easiness assumptions.
\paragraph{Theorem 45.} $\Sigma_2^{\mathrm{QL}}\times\mathrm{QLSamp}\subseteq\mathrm{AvgBPQL}\Longrightarrow\mathrm{AMTIME}[2^{O(\sqrt{n\log n})}]=\mathrm{BPTIME}[2^{O(\sqrt{n\log n})}]$.
It suffices to prove the inclusion $\mathrm{AMTIME}[2^{O(\sqrt{n\log n})}]\subseteq\mathrm{BPTIME}[2^{O(\sqrt{n\log n})}]$ (as the other inclusion is obvious). We will need the following analog of Lemma 31. The difference is that the ``blow-up'' of the time bound in $\mathrm{pK}^t$ here is linear in $t$, while in Lemma 31, this blow-up is polynomial in $t$.
\paragraph{Lemma 46 (Fine-Grained Computational Depth).} Let $x\in\{0,1\}^n$ and $r\in\{0,1\}^*$ be arbitrary strings, $f$ any function, and $\tau_n(t):=t\cdot\operatorname{poly}(n)$. There exists a constant $a>0$ and a time bound $t$ such that $f(n)\leq t\leq f(n)\cdot2^{a\sqrt{n\log n}}$ and
\[\mathrm{pK}^t(x\mid r)-\mathrm{pK}^{\tau_n(t)}(x\mid r)\leq O(\sqrt{n\log n}).\]
\paragraph{Proof.} Fix $x\in\{0,1\}^n$. Let $\tau_n(t):=t\cdot n^c$. Define $t_0:=f(n)$, and for $i\geq1$, $t_i:=\tau_n(t_{i-1})$. Consider the following telescoping sum:
\begin{align*}
\mathrm{pK}^{t_0}(x\mid r)-\mathrm{pK}^{t_I}(x\mid r)
&=(\mathrm{pK}^{t_0}(x\mid r)-\mathrm{pK}^{t_1}(x\mid r))\\
&\quad+(\mathrm{pK}^{t_1}(x\mid r)-\mathrm{pK}^{t_2}(x\mid r))\\
&\quad+\cdots+(\mathrm{pK}^{t_{I-1}}(x\mid r)-\mathrm{pK}^{t_I}(x\mid r)).
\end{align*}
Note that $\mathrm{pK}^{f(n)}(x\mid r)\leq n+d$, for some universal constant $d\geq0$; hence, the above sum is at most $n+d$. By averaging, there is some index $1\leq i^*\leq I$ such that
\[\mathrm{pK}^{t_{i^*-1}}(x\mid r)-\mathrm{pK}^{t_{i^*}(n)}(x\mid r)\leq(n+d)/I,\]
which is at most $O(\sqrt{n\log n})$ for $I:=\sqrt{n/\log n}$. Finally, by induction, for every $i\geq0$, $t_i\leq f(n)\cdot2^{c\cdot i\cdot\log n}$. The bound $t_{i^*-1}\leq f(n)\cdot2^{O(\sqrt{n\cdot\log n})}$ follows.\hfill$\triangleleft$
We will also need the following fine-grained version of Lemma 26.
\paragraph{Lemma 47 (Fine-Grained Symmetry of Information for $\mathrm{pK}^t$).} If $\mathrm{NQL}\times\mathrm{QLSamp}\subseteq\mathrm{AvgBPQL}$, then there exists a polynomial $p$ such that for all sufficiently large $n,m\in\mathbb N$, all $t\in\mathbb N$ such that $t\geq2^{(n\cdot m)^\varepsilon}$ for some $\varepsilon>0$, and all time-constructible $0\leq\tau(n,m)\leq t$, the following holds: for every $x\in\{0,1\}^n$ and every family $\{y_r\in\{0,1\}^m\mid r\in\{0,1\}^{\tau(n,m)}\}$, with probability at least $9/10$ over $r\in\{0,1\}^{\tau(n,m)}$,
\[\mathrm{pK}^t(x,y_r\mid r)>\mathrm{pK}^{\widetilde O(t)}(x\mid r)+\mathrm{pK}^{\widetilde O(t)}(y_r\mid x,r)-\log p(t).\]
\paragraph{Proof Sketch.} The proof is adapted from that of Lemma 26. Define
\begin{multline*}
L:=\{(u,v,w,w',1^t,1^s,1^n,1^m)\mid\exists M\in\{0,1\}^s,M(w,w')\text{ prints }uv\text{ within }2t\text{ steps, where}\\
|w'|=\tau(n,m)\leq t,|w|=2t,\text{ and }s=|u|+|v|-10\}.
\end{multline*}
A distribution $D$ will be given parameters $\langle nk+k,mk'+k',4t,s,n,m\rangle=:N$. It will be defined to randomly sample $u\sim U_{nk+k}$, $v\sim U_{mk'+k'}$, $w\sim U_{2t}$, and $w'\sim U_{\tau(n,m)}$, and then output $(u,v,w,w',1^t,1^s,1^n,1^m)$.
The key observation is that $L$ can be solved in NQL, so using our fine-grained average-case easiness assumption, we obtain a randomized heuristic algorithm $B$ that solves the distributional problem $(L,D)$ in quasilinear time. In particular, since $t\geq2^{(n\cdot m)^\varepsilon}$ for some $\varepsilon>0$, $N\leq\widetilde O(t)$. This ensures that $B$ runs in time $\widetilde O(t)$ (see Proposition 16 and Definition 5).
The rest of the proof follows that of Lemma 26, except when $B$ is used later on as a distinguisher, where we invoke the reconstruction lemma (Lemma 22), we get time bounds of $\widetilde O(t)$ instead of $\operatorname{poly}(t)$.\hfill$\triangleleft$
Using similar ideas, we can get an analog of Lemma 33 for a special case of $\ell=1$. We will use the following definitions, similar to Definition 32, which respectively allow for larger time-bounds and impose a separate length restriction on witnesses.
\paragraph{Definition 48 ($\mathrm{BP}_\delta\circ\mathrm{NTIME}$ and $\mathrm{BP}_\delta\circ\mathrm{NTIMEGUESS}$).} A language $L$ is in the class $\mathrm{BP}_\delta\circ\mathrm{NTIME}[\tau]$ for a time bound $\tau$ if there is a formula
\[\varphi(x)=\mathrm{BP}_r\ \exists y\ R(x,y,r)\]
for a predicate $R$ running in time $\tau$ and variables $y,r$ of length at most $\tau$, such that, for all $x\in\{0,1\}^n$,
\begin{align*}
x\in L&\Longrightarrow\Pr_r[\exists y\ R(x,y,r)]\geq\delta,\text{ and}\\
x\notin L&\Longrightarrow\Pr_r[\forall y\ \neg R(x,y,r)]\geq\delta.
\end{align*}
The class $\mathrm{BP}_\delta\circ\mathrm{NTIMEGUESS}[\tau,\ell]$ is defined the same way as above, with an extra condition that $|y|\leq\ell$ (but the randomness $r$ may still be of length at most $\tau$).
\paragraph{Lemma 49 (Fine-Grained Oracle Elimination).} Suppose $\mathrm{NQL}\times\mathrm{QLSamp}\subseteq\mathrm{AvgBPQL}$. Let $\tau(n)=2^{O(\sqrt{n\log n})}$ be some time-constructible function and $c>0$ a constant. For an arbitrary $L\in\mathrm{BP}_\delta\circ\mathrm{NTIMEGUESS}[\tau(n),n^c]$, let $x\in L_n$ be sufficiently large, let $r\in\{0,1\}^{\tau(n)}$ be any random string under the BP quantifier such that there exists a witness for $x\in L$, and let $y_{x,r}\in\{0,1\}^{n^c}$ be the lexicographically first such $L$-witness for $x$ and randomness $r$.
There exist polynomials $q$ and $q_0$ such that, for all sufficiently large $n\in\mathbb N$, all $x\in L_n$, and all $t\geq q_0(\tau(n))$, we have with probability at least $\delta-1/10$ over $r\in\{0,1\}^{\tau(n)}$ that an $L$-witness $y_{x,r}$ exists, and
\[\mathrm{pK}^{\widetilde O(t)}(x,y_{x,r}\mid r)\leq\mathrm{pK}^t(x\mid r)+\log q(t).\]
\paragraph{Proof Sketch.} The proof is adapted from that of Lemma 33. Define
\begin{multline*}
L':=\{(\mathrm{DP}_k(x,y;z),w,w',1^t,1^k,1^s,1^n)\mid\exists M\in\{0,1\}^s,\text{ such that }M^{\mathrm{SAT}}(w,w')\text{ prints}\\
(x,y)\in\{0,1\}^{n+n^c}\text{ within }2t\text{ steps, where }|w'|=\tau(n)\leq t,|w|=2t,\text{ and }s=k-10\}.
\end{multline*}
Since SAT is NQL-complete under many-one QL reduction, we have, for $x\in\{0,1\}^n$, $y\in\{0,1\}^{n^c}$, $w\in\{0,1\}^{2t}$, and $w'\in\{0,1\}^{\tau(n)}$ that $L'\in\Sigma_2^{\mathrm{QL}}$.
Again, using our fine-grained average-case easiness assumption, we get a randomized heuristic algorithm $B$ that solves $L'$ in time $\widetilde O(t)$. Also note that given $x$ and good randomness $r\in\{0,1\}^{\tau(n)}$, we can compute $y_{x,r}$ via search-to-decision reduction using a SAT oracle, which takes time $\operatorname{poly}(\tau)=:q_0(\tau)$. Therefore, similar to Equation (19) in the original proof, we get
\[\mathrm{pK}^{2t,\mathrm{SAT}}(x,y_{x,r}\mid r)\leq\mathrm{pK}^t(x\mid r)+d\log t\tag{30}\]
for some constant $d>0$, provided $t\geq q_0(\tau)$. The rest of the proof is the same as that of Lemma 33, except when $B$ is used later in the proof as a distinguisher, where we invoke the reconstruction lemma (Lemma 22), we get a time bound of $\widetilde O(t)$ instead of $\operatorname{poly}(t)$.\hfill$\triangleleft$
We will also use the following simple generalization of (non-fine-grained) oracle elimination, Lemma 33, to higher time bounds.
\paragraph{Lemma 50.} Suppose $\mathrm{Dist}\Sigma_2^{\mathrm P}\subseteq\mathrm{AvgBPP}$. For an arbitrary $L\in\mathrm{NTIME}[\tau(n)]$, for some $\tau(n)=2^{O(\sqrt{n\log n})}$, let $x\in L_n$ be sufficiently large, and let $y_x\in\{0,1\}^{\tau(n)}$ be the lexicographically first $L$-witness for $x$. There exist polynomials $q$ and $q_0$ such that, for all $t\geq q_0(\tau(n))$, we have
\[\mathrm{pK}^{q(t)}(x,y_x)\leq\mathrm{pK}^t(x)+\log q(t).\]
\paragraph{Proof Sketch.} The proof follows that of the ``Moreover'' statement of Lemma 33, except the search-to-decision computation of $y_x$, where $|y_x|=\tau(n)$, now takes time at least $q_0(\tau(n))$ for some polynomial $q_0$, yielding
\[\mathrm{pK}^{2t,\mathrm{NP}}(x,y_{x,r}\mid r)\leq\mathrm{pK}^t(x\mid r)+d\log n\]
for $t\geq q_0(\tau(n))$ in Equation (19).\hfill$\triangleleft$
We are now ready to show Theorem 45.
\paragraph{Proof of Theorem 45.} Assume $\Sigma_2^{\mathrm{QL}}\times\mathrm{QLSamp}\subseteq\mathrm{AvgBPQL}$. Let $\tau=2^{a\sqrt{n\log n}}$ for some constant $a>0$, and let $L\in\mathrm{AMTIME}[\tau]$ be arbitrary. That is, for every $x\in\{0,1\}^n$,
\begin{align*}
x\in L&\Longrightarrow\mathrm{BP}_{(1-2^{-n})}r\ \exists y\ V(x,y,r),\\
x\notin L&\Longrightarrow\mathrm{BP}_{(1-2^{-n})}r\ \forall y\ \neg V(x,y,r),
\end{align*}
for a polytime verifier $V(x,y,r)$, where $|r|,|y|=\tau$. Let $x\in L_n$ and good randomness $r\in\{0,1\}^\tau$ be given, with lexicographically first witness $y_{x,r}\in\{0,1\}^\tau$.
Note that $\Sigma_2^{\mathrm{QL}}\times\mathrm{QLSamp}\subseteq\mathrm{AvgBPQL}$ implies $\mathrm{Dist}\Sigma_2^{\mathrm P}\subseteq\mathrm{AvgBPP}$ by a simple padding argument. Combining Lemma 50 and Lemma 26 as in the proof of Lemma 34, for some constant $b>0$, with probability at least $1-2^{-n}-1/5$ over $r\in\{0,1\}^\tau$,
\[\mathrm{pK}^{\tau^b}(y_x\mid x,r)\leq(\mathrm{pK}^\tau(x\mid r)-\mathrm{pK}^{\tau^b}(x\mid r))+O(\log\tau),\tag{31}\]
from which it follows that $\mathrm{pK}^{\tau^b}(y_x\mid x,r)\leq2n$ by the trivial upper bound on $\mathrm{pK}^\tau(x\mid r)$. In other words, by definition of pK,
\[\Pr_{r'\sim\{0,1\}^{\tau^b}}[\exists M\in\{0,1\}^{2n},M(x,r,r')\text{ outputs }y_x\text{ within time }\tau^b]\geq2/3.\tag{32}\]
Define a verifier $V'(x,z,r,r')$, which does the following:
\begin{itemize}
\item Simulate the machine described by $z$ on input $(x,r,r')$ for $\tau^b$ steps to produce some output $y_z$, and then return the output of $V(x,y_z,r)$.
\end{itemize}
Note that $V'$ runs in time $\tau^c$ for some constant $c>0$. For $x\in L$, there exists $z\in\{0,1\}^{2n}$ such that $V'(x,z,r,r')$ so long as Equation (31) and Equation (32) both hold, which occurs with probability at least $(2/3)(1-2^{-n}-1/5)=:\delta>1/2$ over $r,r'$. For $x\notin L$, $V'(x,z,r,r')$ rejects unless there is an $L$-witness $y_{x,r}$ under the original verifier $V$, which occurs with probability at most $2^{-n}$. This leads to the following characterization of $L$:
\begin{align*}
x\in L&\Longrightarrow\mathrm{BP}_{(\delta)}r,r'\ \exists z\ V'(x,z,r.r'),\\
x\notin L&\Longrightarrow\mathrm{BP}_{(1-2^{-n})}r,r'\ \forall z\ \neg V'(x,z,r,r'),
\end{align*}
where $|r|=\tau$, $|r'|=\tau^b$, and $|z|=2n$, which implies that $L\in\mathrm{BP}_\delta\circ\mathrm{NTIMEGUESS}[\tau^c,2n]$.
Combining Lemma 47 and Lemma 49 as in the proof of Lemma 34, we have that with probability at least $\delta-1/5>1/4$ over $(r,r')\in\{0,1\}^{\tau+\tau^b}$ that for all $t\geq p'_0(\tau^c)$, for some polynomials $p'$ and $p'_0$,
\[\mathrm{pK}^{t\cdot p'(n)}(z_{x,r,r'}\mid x,r,r')\leq(\mathrm{pK}^t(x\mid r,r')-\mathrm{pK}^{t\cdot p'(n)}(x\mid r,r'))+O(\log t)\tag{33}\]
where $z_{x,r,r'}$ is such that $V'(x,z_{x,r,r'},r,r')$ accepts. Lemma 46 then implies the existence of some $t\in\mathbb N$ such that $p'_0(\tau^c)\leq t\leq2^{O(\sqrt{n\log n})}$ and $\mathrm{pK}^t(x\mid r,r')-\mathrm{pK}^{t\cdot p'(n)}(x\mid r,r')\leq O(\sqrt{n\log n})$. So by Equation (33), for some constant $d>0$,
\[\mathrm{pK}^{\tau^d}(z_{x,r}\mid x,r,r')\leq d\sqrt{n\log n}.\]
This suggests the following algorithm to decide $L$:
\begin{itemize}
\item Given $x\in\{0,1\}^n$, sample uniformly random strings $r\sim\{0,1\}^\tau$, $r'\in\{0,1\}^{\tau^b}$, and $w\in\{0,1\}^{\tau^d}$. Then exhaustively search over all $d\sqrt{n\log n}$-bit machines $M$, running each $M(x,r,r',w)$ for $\tau^d$ steps to produce a candidate witness $z$, accepting if $V'(x,z,r,r')$ accepts.
\end{itemize}
The runtime $\operatorname{poly}(\tau)=2^{O(\sqrt{n\log n})}$ is clear. For correctness, for every $x\in L_n$, there is at least a $1/4$ fraction of good random strings $r,r'$ such that Equation (33) holds, in which case a witness $z_{x,r,r'}$ will be found by our algorithm with probability at least $2/3$ over $w$ (by definition of pK). So, our algorithm accepts $x\in L_n$ with probability at least $(1/4)(2/3)=1/6$. On the other hand, for $x\notin L$, our algorithm will accept with probability at most $2^{-n}$, as this upper-bounds the probability of $r$ being good for the original verifier $V$. Thus, we have $L\in\mathrm{BPTIME}[2^{O(\sqrt{n\log n})}]$ via standard success amplification.\hfill$\triangleleft$
A proof of the analogous statement for the class $\mathrm{UTIME}[2^{O(\sqrt{n\log n})}]$, Item 1 of Theorem 2,
\[\mathrm{NQL}\times\mathrm{QLSamp}\subseteq\mathrm{AvgBPQL}\Longrightarrow\mathrm{UTIME}[2^{O(\sqrt{n\log n})}]\subseteq\mathrm{RTIME}[2^{O(\sqrt{n\log n})}],\]
is presented in Appendix B. Though its proof requires some additional ideas, it does not need to use an extremely efficient PRG as the one developed in [6].
